\scanpage{67}
Notons le \uline{th. de dualité absolue} : % 1
Soient $M$ et $M^*$ deux \inserted{(complexes de)} motifs sur % 2
$X$ (de type fini sur $\mathbb{Z}$), qui sont % 3
\textquotedblleft duals\textquotedblright{} pour le \deleted{\unsure{\textup{[}}}complexe de motifs % 4
dualisant\deleted{\unsure{s}}\deleted{\unsure{\textup{]}}} canonique -- [NB. Si $X$ % 5
est régulier \inserted{partout} de dim $n$, alors % 6
c'est la dualité à valeurs \deleted{dans $\unsure{\mathbb{Q}(n)[\,]| \mathbb{Q}(n)[2n]}$ shift} \inserted{$\mathbb{Q}(n)[2n+1]$}\margin{shift}] % 7

alors $\mathbb{R}\Gamma_X(M)$ et $\mathbb{R}_!\Gamma_X(M^*)$\margin{\textquotedblleft support cpt sur $\operatorname{Spec}\mathbb{Z}$\textquotedblright} sont % 8
duals. Ici [Comme les coefficients sont % 9
ess. $\mathbb{Q}$\deleted{\unsure{s|\textup{\}}}}, donc il n'y a pas de $\operatorname{Tor}$] % 10
ceci revient aussi à la formule (cas régulier) % 11
\[
\boxed{\begin{aligned}
H^i(X,M)&\quad\text{dual de}\quad\operatorname{Ext}^{2n+1-i}_!(X;M,\mathbb{Q}(n))\\ % 12
H^i_!(X,M)&\quad\text{dual de}\quad\operatorname{Ext}^{2n+1-i}(X;M,\mathbb{Q}(n)) % 13
\end{aligned}}
\]
Si $X$ propre sur $\operatorname{Spec}\mathbb{Z}$, inutile de mettre % 14
des $!$. % 15

Moyennant le \unsure{thm|th.} de dualité \uline{relative} % 16
pour motifs, cet énoncé revient au % 17
% archive sheet number: 66
